\subsection{Native polynomial Proj and the normalized graded seminorm model}
\label{subsec:native-normalized-projective-comparison}

The comparison is constructed from actual prime supports, homogeneous
localizations and the true projective structure sheaf. The coefficient
ring and its prescribed real seminorm are arbitrary. No chart certificate,
residue valuation or global analytification identification is assumed.
The native residue field below is the quotient of the genuine degree-zero
projective stalk. It is not identified with the generally larger residue
field of the affine cone seminorm.

\begin{definition}[The genuine support and projective stalk seminorm]
\label{projective:normalized-native-support-stalk}
\lean{EnvelopeProperty.Berkovich.NormalizedProjectiveSpectrum.toRelevantGradedSpectrum,
  EnvelopeProperty.Berkovich.NormalizedProjectiveSpectrum.continuous_toRelevantGradedSpectrum,
  EnvelopeProperty.Berkovich.NormalizedProjectiveSpectrum.nativeSupport,
  EnvelopeProperty.Berkovich.NormalizedProjectiveSpectrum.continuous_nativeSupport,
  EnvelopeProperty.Berkovich.NormalizedProjectiveSpectrum.nativeSupport_toIdeal,
  EnvelopeProperty.Berkovich.NormalizedProjectiveSpectrum.nativeSupport_mem_coordinateOpen,
  EnvelopeProperty.Berkovich.NormalizedProjectiveSpectrum.localValuation,
  EnvelopeProperty.Berkovich.NormalizedProjectiveSpectrum.localValuation_supp,
  EnvelopeProperty.Berkovich.NormalizedProjectiveSpectrum.homogeneousLocalValuation,
  EnvelopeProperty.Berkovich.NormalizedProjectiveSpectrum.homogeneousLocalValuation_supp,
  EnvelopeProperty.Berkovich.NormalizedProjectiveSpectrum.nativeStalkValuation,
  EnvelopeProperty.Berkovich.NormalizedProjectiveSpectrum.nativeStalkValuation_supp,
  EnvelopeProperty.ProjBerk.iSup_nativeCoordinateOpen_eq_top}
\leanok
\uses{projective:general-normalized-spectrum,spectrum:graded-support,spectrum:relevant-graded-support,projective:native-scheme-coordinate-sections}
A normalized graded polynomial seminorm defines its actual homogeneous
prime kernel, avoiding the irrelevant ideal, and therefore a point of the
actual polynomial Proj scheme. Extend the original seminorm to its ordinary
localization at that prime, restrict it to the genuine homogeneous degree-
zero localization, and transport it through the true projective stalk
isomorphism. Its support is exactly the actual stalk maximal ideal.
\end{definition}
\begin{proof}
The actual homogeneous projection inequality makes the valuation kernel
homogeneous. Some normalized coordinate has value one; that coordinate
belongs to the irrelevant ideal but not the kernel, proving relevance.
Every denominator outside the kernel has nonzero seminorm, so actual
valuation localization is available. In the ordinary localization, a
fraction has zero value precisely when its numerator belongs to the
original prime. The true localization maximal-ideal criterion therefore
identifies the localized support with its maximal ideal. A homogeneous
fraction is a unit exactly when its image in the ordinary local ring is,
so the restricted homogeneous-localization support is also its maximal
ideal. The genuine projective stalk ring isomorphism preserves units,
proving the stated support identity in the actual stalk.
\end{proof}

\begin{theorem}[Actual native residue values and every coordinate section]
\label{projective:normalized-native-residue-values}
\lean{EnvelopeProperty.Berkovich.NormalizedProjectiveSpectrum.nativePoint,
  EnvelopeProperty.Berkovich.NormalizedProjectiveSpectrum.nativePoint_point,
  EnvelopeProperty.Berkovich.NormalizedProjectiveSpectrum.nativePoint_residue,
  EnvelopeProperty.Berkovich.NormalizedProjectiveSpectrum.homogeneousLocalValuation_mapId_chart,
  EnvelopeProperty.Berkovich.NormalizedProjectiveSpectrum.nativePoint_eval_homogeneousChart,
  EnvelopeProperty.Berkovich.NormalizedProjectiveSpectrum.nativePoint_eval_coordinateChart,
  EnvelopeProperty.ProjBerk.nativeAbsoluteCoordinateHomeomorph,
  EnvelopeProperty.ProjBerk.nativeAbsoluteCoordinateHomeomorph_val_apply}
\leanok
\uses{projective:normalized-native-support-stalk,projective:native-scheme-coordinate-sections,spectrum:localization-extension,scheme:affine-chart-section-values}
The actual stalk valuation descends to the actual residue field of that
projective point and defines a genuine native valued scheme point. In
any coordinate chart, its value on every affine polynomial equals the
whole ratio seminorm obtained from the original normalized point.
Its underlying projective scheme support is the original homogeneous
prime kernel.
\end{theorem}
\begin{proof}
The proved maximal-ideal support supplies genuine quotient descent to the
actual local residue field. The descended value on every residue of a
stalk element equals the original stalk value.

A coordinate is nonzero precisely when this true projective support lies
in that coordinate open. The true affine-open section homeomorphism and the actual coordinate
section-ring isomorphism also give an absolute native coordinate chart,
preserving every coordinate polynomial section. Its injectivity identifies
whole native points with matching affine polynomial seminorms.

The actual section/germ comparison identifies
each homogeneous fraction section with its image in the true homogeneous
local ring at this support. Inclusion into the full ordinary localization
preserves that entire fraction. Valuation localization uniqueness identifies
this with the genuine principal-localization extension of the original
seminorm. The actual ratio homomorphism therefore gives the stated value
identity on every affine polynomial, rather than only on generators.
\end{proof}

\begin{theorem}[Actual prescribed-base native projective homeomorphism]
\label{projective:normalized-native-homeomorphism}
\lean{EnvelopeProperty.Berkovich.RelativeNormalizedProjectiveSpectrum.nativePoint_eq_chart_inverse,
  EnvelopeProperty.Berkovich.RelativeNormalizedProjectiveSpectrum.nativePoint_mem_relativeLocus,
  EnvelopeProperty.Berkovich.RelativeNormalizedProjectiveSpectrum.nativePoint,
  EnvelopeProperty.Berkovich.RelativeNormalizedProjectiveSpectrum.nativePoint_val,
  EnvelopeProperty.Berkovich.RelativeNormalizedProjectiveSpectrum.nativePoint_onChart,
  EnvelopeProperty.Berkovich.RelativeNormalizedProjectiveSpectrum.nativePoint_injective,
  EnvelopeProperty.Berkovich.RelativeNormalizedProjectiveSpectrum.nativePoint_surjective,
  EnvelopeProperty.Berkovich.RelativeNormalizedProjectiveSpectrum.continuous_nativePoint,
  EnvelopeProperty.Berkovich.RelativeNormalizedProjectiveSpectrum.nativeHomeomorph,
  EnvelopeProperty.Berkovich.RelativeNormalizedProjectiveSpectrum.nativeHomeomorph_apply,
  EnvelopeProperty.Berkovich.RelativeNormalizedProjectiveSpectrum.nativeHomeomorph_point,
  EnvelopeProperty.Berkovich.RelativeNormalizedProjectiveSpectrum.nativeHomeomorph_support_toIdeal,
  EnvelopeProperty.Berkovich.RelativeNormalizedProjectiveSpectrum.nativeHomeomorph_val_apply,
  EnvelopeProperty.Berkovich.RelativeNormalizedProjectiveSpectrum.nativeHomeomorph_chart,
  EnvelopeProperty.Berkovich.RelativeNormalizedProjectiveSpectrum.nativeHomeomorph_symm_onChart,
  EnvelopeProperty.Berkovich.RelativeNormalizedProjectiveSpectrum.nativeHomeomorph_eval_form_of_val_X_eq_one}
\leanok
\uses{projective:normalized-native-residue-values,projective:general-relative-normalized-charts,projective:native-relative-coordinate-chart,scheme:separated-native-hausdorff}
The normalized graded spectrum with prescribed coefficient seminorm $v$
is homeomorphic to the actual native valued-point space of polynomial
Proj over the whole base point $v$. This homeomorphism preserves the
actual projective prime support and every genuine coordinate-chart
polynomial evaluation.
\end{theorem}
\begin{proof}
The preceding canonical residue construction supplies the point map.
On a coordinate open, its whole evaluation identity identifies it with
the composite of the genuine normalized relative chart and the inverse
actual native relative chart. The proved coefficient structure diagram
shows that it lies above the entire original prescribed base point.

The actual coordinate opens cover Proj because the true irrelevant ideal
is the ideal generated by all coordinate variables. Given any native
point, choose such a chart, use its genuine affine seminorm, and apply
the actual normalized inverse Gauss chart. The whole value identity and
native chart injectivity recover that entire native point, proving
surjectivity. For injectivity, equal native points have equal support;
choose a nonzero coordinate for one normalized point, hence for both.
The genuine native chart and normalized chart inverse recover the same
whole normalized point.

Each normalized coordinate open is open, and the just-proved chart
identity expresses the map there as a composite of genuine continuous
chart maps. These opens cover, proving global continuity. The normalized
prescribed-base space is genuinely compact for arbitrary $v$. Polynomial
Proj is genuinely separated; the actual native topology of a separated
scheme is Hausdorff. A continuous bijection from this compact space to
that Hausdorff native space is a homeomorphism. The homeomorphism keeps the genuine homogeneous prime ideal equal to
the original valuation kernel and commutes with every entire relative
coordinate chart; its inverse is the true Gauss inverse on any such
chart. When the selected coordinate has value one, actual section
evaluation of a degree-$d$ form $F/X_j^d$ equals the original normalized
value of $F$, since its actual fraction denominator has value one.

No boundedness of $v$,
field, completeness, noetherian or coherent-sheaf hypothesis enters the
argument. This identifies point spaces and topology; an analytic structure
sheaf is a separate construction.
\end{proof}
